\section{符号说明}\label{sec:symbols}

本文统一使用表~\ref{tab:symbols} 所列符号；仅在局部使用的中间量，均在首次出现处说明。

\begin{table}[H]
\centering
\caption{主要符号及其含义}
\label{tab:symbols}
\small
\renewcommand{\arraystretch}{1.12}
\setlength{\tabcolsep}{6pt}
\begin{tabular}{
  >{\centering\arraybackslash}p{2.8cm}
  p{8.8cm}
  >{\centering\arraybackslash}p{2.8cm}}
\toprule
符号 & 含义 & 单位或取值 \\
\midrule
\(i,j\) & 催化剂组合编号、组合内按温度升序排列的观测序号 & 正整数 \\
\(n_i\) & 催化剂组合 \(i\) 的温度观测数 & 个 \\
\(T,t\) & 反应温度、实验时间 & ℃，min \\
\(X_i(T)\) & 组合 \(i\) 在温度 \(T\) 下的乙醇转化率 & \% \\
\(S_i(T)\) & 组合 \(i\) 在温度 \(T\) 下的 C4 烯烃选择性 & \% \\
\(Y_i(T)\) & 组合 \(i\) 在温度 \(T\) 下的 C4 烯烃收率 & \% \\
\(X_{ij},S_{ij},Y_{ij}\) & 上述三项指标在组合 \(i\) 第 \(j\) 个温度观测处的离散记号 & \% \\
\(\Delta X_{ij},\Delta S_{ij}\) & 相邻温度观测之间的转化率差和选择性差 & 百分点 \\
\(Z_{ij},\widehat Z_{ij}\) & 考察指标的观测值及其线性拟合值 & \% \\
\(a_i,b_i\) & 组合 \(i\) 的线性回归截距和斜率 & \%，\%/℃ \\
\(e_{ij},R_i^2\) & 线性回归残差和决定系数 & 百分点，无量纲 \\
\(v_j\) & 相邻时刻之间折算到每 10 min 的指标变化量 & 百分点 \\
\(P_k(t)\) & 时刻 \(t\) 第 \(k\) 类产物的选择性 & \% \\
\(r_A(t),r_H(t)\) & 乙醛、脂肪醇相对于 C4 烯烃的选择性份额比 & 无量纲 \\
\(y_{ij}\) & 共同温度分析中组合 \(i\) 在温度 \(j\) 下的考察指标 & \% \\
\(\mu\) & 相加模型中的总体均值 & \% \\
\(\alpha_i\) & 扣除温度总体差异后的组合偏差 & 百分点 \\
\(\beta_j\) & 温度平均偏差 & 百分点 \\
\(r_{ij}\) & 简单相加未解释变化 & 百分点 \\
\(D_{0}^{\mathrm{all}},D_{0}^{\mathrm{low}}\) & 第三问中的全部实测候选域和严格低温实测候选域 & 集合 \\
\(c_e,T_e\) & 第 \(e\) 次新增实验的催化剂组合和反应温度 & 组合，℃ \\
\(X_e,S_e,Y_e\) & 第 \(e\) 次新增实验测得的转化率、选择性和收率 & \% \\
\bottomrule
\end{tabular}
\end{table}

除特别说明外，本文用百分数表示比例，用``百分点''表示两个百分数的绝对差。例如，
选择性由 30\% 增至 35\% 时，其变化量为 5 个百分点。
